% IEEE/ACM Conference Paper Template
% 双栏格式，适用于计算机安全/网络会议论文

\documentclass[10pt,conference]{IEEEtran}

% 基础包
% 注意：如果使用 pdflatex 编译，保留下面两行；如果使用 xelatex，注释掉
\usepackage[utf8]{inputenc}
\usepackage[T1]{fontenc}

\usepackage{cite}
\usepackage{amsmath,amssymb,amsfonts}
\usepackage{algorithmic}
\usepackage{algorithm}
\usepackage{graphicx}
\usepackage{textcomp}
\usepackage{xcolor}
\usepackage{hyperref}
\usepackage{url}
\usepackage{booktabs}
\usepackage{multirow}
\usepackage{subcaption}

% 页面设置
\usepackage[letterpaper, margin=0.75in]{geometry}

% 超链接设置
\hypersetup{
    colorlinks=true,
    linkcolor=black,
    citecolor=blue,
    filecolor=magenta,
    urlcolor=blue,
}

% 定义一些常用命令
\def\BibTeX{{\rm B\kern-.05em{\sc i\kern-.025em b}\kern-.08em
    T\kern-.1667em\lower.7ex\hbox{E}\kern-.125emX}}

\begin{document}

% ========================================
% 标题部分
% ========================================
\title{Your Paper Title Here: \\
Make It Descriptive and Clear}

% 作者信息（匿名审稿版本可以注释掉）
\author{
\IEEEauthorblockN{First Author\IEEEauthorrefmark{1},
Second Author\IEEEauthorrefmark{2},
Third Author\IEEEauthorrefmark{1}}
\IEEEauthorblockA{\IEEEauthorrefmark{1}First Institution\\
Department Name\\
City, Country\\
Email: \{first.author, third.author\}@institution.edu}
\IEEEauthorblockA{\IEEEauthorrefmark{2}Second Institution\\
Department Name\\
City, Country\\
Email: second.author@institution.edu}
}

\maketitle

% ========================================
% 摘要
% ========================================
\begin{abstract}
This is the abstract of your paper. The abstract should briefly summarize the contents of the paper in 150--250 words. It should clearly state the problem you are addressing, your approach, key findings, and main contributions. The abstract is typically written last, after you have completed the paper.
\end{abstract}

% ========================================
% 关键词（某些会议需要）
% ========================================
% \begin{IEEEkeywords}
% keyword1, keyword2, keyword3, keyword4
% \end{IEEEkeywords}

% ========================================
% I. 引言
% ========================================
\section{Introduction}
\label{sec:introduction}

The introduction should motivate the problem and clearly state your contributions. Start with the broader context, then narrow down to the specific problem you are addressing.

\subsection{Motivation}
Explain why this problem is important and worth solving. Provide context and background that helps readers understand the significance of your work.

\subsection{Problem Statement}
Clearly define the problem you are addressing. What are the gaps in existing solutions? What challenges need to be overcome?

\subsection{Contributions}
List your main contributions in a clear, numbered list:
\begin{itemize}
    \item \textbf{First contribution:} Brief description of the first major contribution.
    \item \textbf{Second contribution:} Brief description of the second major contribution.
    \item \textbf{Third contribution:} Brief description of the third major contribution.
\end{itemize}

\subsection{Paper Organization}
The rest of this paper is organized as follows. Section~\ref{sec:background} provides background and related work. Section~\ref{sec:methodology} describes our methodology. Section~\ref{sec:evaluation} presents our evaluation. Section~\ref{sec:discussion} discusses the implications and limitations. Finally, Section~\ref{sec:conclusion} concludes the paper.

% ========================================
% II. 背景与相关工作
% ========================================
\section{Background and Related Work}
\label{sec:background}

\subsection{Background}
Provide the necessary background information that readers need to understand your work. Define key terms and concepts.

\subsection{Related Work}
Discuss prior work related to your research. Compare and contrast your approach with existing solutions. Cite relevant papers~\cite{example1, example2}.

% ========================================
% III. 方法论/系统设计
% ========================================
\section{Methodology}
\label{sec:methodology}

Describe your approach, system design, or methodology in detail. This section should be substantial and may be divided into multiple subsections.

\subsection{Overview}
Provide a high-level overview of your approach. Consider including a figure that illustrates the overall architecture or workflow.

% 取消注释下面的代码来插入图片（需要先将图片放在 figures/ 目录中）
% \begin{figure}[t]
%     \centering
%     \includegraphics[width=0.9\columnwidth]{figures/system_overview.pdf}
%     \caption{System overview showing the main components and their interactions.}
%     \label{fig:overview}
% \end{figure}

\subsection{Component A}
Describe the first major component of your system or methodology.

\subsection{Component B}
Describe the second major component.

\subsection{Component C}
Describe additional components as needed.

% ========================================
% IV. 实现
% ========================================
\section{Implementation}
\label{sec:implementation}

Describe implementation details, including:
\begin{itemize}
    \item Programming languages and frameworks used
    \item Hardware/software requirements
    \item Key implementation challenges and solutions
\end{itemize}

% ========================================
% V. 评估/实验
% ========================================
\section{Evaluation}
\label{sec:evaluation}

Present your experimental results and evaluation.

\subsection{Experimental Setup}
Describe your experimental methodology:
\begin{itemize}
    \item Dataset description
    \item Evaluation metrics
    \item Baseline comparisons
    \item Experimental environment
\end{itemize}

\subsection{Results}
Present your main results. Use tables and figures to clearly communicate your findings.

\begin{table}[t]
\centering
\caption{Performance comparison of different approaches}
\label{tab:performance}
\begin{tabular}{@{}lccc@{}}
\toprule
Method & Metric 1 & Metric 2 & Metric 3 \\
\midrule
Baseline 1 & 85.2\% & 12.3ms & 95.1\% \\
Baseline 2 & 87.5\% & 10.8ms & 96.2\% \\
\textbf{Our Approach} & \textbf{92.3\%} & \textbf{8.5ms} & \textbf{98.7\%} \\
\bottomrule
\end{tabular}
\end{table}

\subsection{Analysis}
Analyze and interpret your results. What do the results tell us? How do they support your claims?

% ========================================
% VI. 讨论
% ========================================
\section{Discussion}
\label{sec:discussion}

\subsection{Implications}
Discuss the broader implications of your findings.

\subsection{Limitations}
Be honest about the limitations of your approach. What are the constraints or assumptions? What scenarios does your solution not handle well?

\subsection{Future Work}
Suggest directions for future research building on your work.

% ========================================
% VII. 对策/安全考量（如适用）
% ========================================
\section{Countermeasures}
\label{sec:countermeasures}

If your paper discusses attacks or vulnerabilities, include a section on potential countermeasures or defenses.

% ========================================
% VIII. 相关工作（可选，如果与背景合并）
% ========================================
% \section{Related Work}
% \label{sec:related}
%
% More detailed discussion of related work, if not covered in Section~\ref{sec:background}.

% ========================================
% IX. 结论
% ========================================
\section{Conclusion}
\label{sec:conclusion}

Summarize your main findings and contributions. Reiterate why your work is important and what impact it may have. Keep this section concise.

% ========================================
% 致谢（可选）
% ========================================
\section*{Acknowledgments}
We thank the anonymous reviewers for their valuable feedback. This work was supported by [funding source].

% ========================================
% 参考文献
% ========================================
\bibliographystyle{IEEEtran}
% 使用 .bib 文件（取消注释下面这行，并确保 references.bib 文件存在）
% \bibliography{references}

% 手动输入参考文献（示例）
\begin{thebibliography}{99}
\bibitem{example1}
A. Author, B. Author, and C. Author,
``Title of the paper,''
in \textit{Proceedings of Conference Name},
2024, pp. 1--10.

\bibitem{example2}
D. Author and E. Author,
``Another paper title,''
\textit{Journal Name},
vol. 10, no. 2, pp. 100--120,
2023.
\end{thebibliography}

\end{document}
